\documentclass{article}
\usepackage{amsmath}
\usepackage{booktabs}
\usepackage{array}
\usepackage{multirow}
\usepackage{xcolor}
\usepackage{colortbl}
\definecolor{headerblue}{RGB}{23, 75, 137}
\definecolor{lightblue}{RGB}{217, 234, 250}
\definecolor{cellgreen}{RGB}{221, 241, 221}
\definecolor{cellyellow}{RGB}{254, 243, 199}
\definecolor{cellred}{RGB}{253, 226, 226}
\begin{document}
\section{Question 1}
\[
\overbrace{a + b + c }^{\text{First Group}} + \underbrace{x+y+z}_{\text{Second Group}}
=
\overbrace{\underbrace{a + b}_{\text{Sum 1}} + c}^{\text{Total}}
\]
\section{Question 2}
\[
A =
\begin{bmatrix}
\begin{array}{c|ccc}
& x_1 & x_2 & x_3 \\
\hline
y_1 & a_{11} & a_{12} & a_{13} \\
y_2 & a_{21} & a_{22} & a_{23} \\
y_3 & a_{31} & a_{32} & a_{33} \\
\end{array}
\end{bmatrix}
\]
\section{Question 3}
\[
f(x) =
\begin{cases}
x^2 + 2x + 1, & x < 0, \\
\sqrt{x} + 1, & 0 \le x < 4, \\
\dfrac{x^2 - 1}{x - 1}, & x \ge 4.
\end{cases}
\]
\section{Question 4}
\[
\begin{bmatrix}
\overbrace{\quad a_{11} \quad a_{12} \quad a_{13}}^{\text{Column Group 1}}
\overbrace{ a_{14} \quad a_{15}}^{\text{Column Group 2}} \\
a_{21} \quad a_{22} \quad a_{23} \quad a_{24} \quad a_{25} \\
a_{31} \quad a_{32} \quad a_{33} \quad a_{34} \quad a_{35}
\end{bmatrix}
\]
\section{Question 5}
\[
F(x) =
\underbrace{
\overbrace{
\sum_{i=1}^{n} \left(a_i x_i^2 + b_i x_i + c_i \right)
}^{\text{Quadratic Contribution}}
+
\overbrace{
\begin{bmatrix}
x_1 & x_2 & \cdots & x_n
\end{bmatrix}
A
\begin{bmatrix}
x_1 \\
x_2 \\
\vdots \\
x_n
\end{bmatrix}
}^{\text{Matrix Contribution}}
}_{\text{Total Energy}}
\]
\[
\text{where }
A =
\begin{bmatrix}
a_{11} & a_{12} & \cdots & a_{1n} \\
a_{21} & a_{22} & \cdots & a_{2n} \\
\vdots & \vdots & \ddots & \vdots \\
a_{n1} & a_{n2} & \cdots & a_{nn}
\end{bmatrix}
\]
\section{Question 6}
\begin{center}
\begin{tabular}{|c|c|c|c|c|}
\hline
\textbf{Roll No.} & \textbf{Name} & \textbf{Department} & \textbf{Semester} & \textbf{CGPA} \\
\hline
101 & Amit  & CSE & 2 & 8.45 \\
102 & Neha  & CSE & 2 & 9.12 \\
103 & Rahul & ECE & 4 & 8.76 \\
104 & Priya & ECE & 4 & 9.01 \\
\hline
\end{tabular}
\end{center}
\section{Question 7}
\begin{center}
\begin{tabular}{|c|c|c|}
\hline
\multicolumn{3}{|c|}{\textbf{Academic Performance}} \\
\hline
\textbf{Student} & \textbf{Midterm} & \textbf{Final} \\
\hline
A & 78 & 89 \\
B & 85 & 92 \\
C & 91 & 94 \\
\hline
\end{tabular}
\end{center}
\noindent\textbf{Que 8} Create the following table.

\begin{center}
\begin{tabular}{|c|l|c|c|}
\hline
\textbf{Department} & \textbf{Course} & \textbf{Credits} & \textbf{Grade} \\
\hline
\multirow{3}{*}{CSE} & Data Structures   & 4 & A  \\
                     & Machine Learning  & 4 & A+ \\
                     & Operating Systems & 3 & B+ \\
\hline
\multirow{2}{*}{ECE} & Digital Electronics & 4 & A  \\
                     & Signals             & 3 & B+ \\
\hline
\end{tabular}
\end{center}

\noindent Use \texttt{\textbackslash multirow} so that CSE and ECE each occupy multiple rows.

\vspace{1.5em}

\noindent\textbf{Que 9} Create the following scientific table.

\begin{center}
\begin{tabular}{|l|c|c|c|c|}
\hline
\rowcolor{headerblue}
\textcolor{white}{\textbf{Model}} & 
\textcolor{white}{\textbf{Parameters}} & 
\textcolor{white}{\textbf{FLOPs}} & 
\textcolor{white}{\textbf{Accuracy}} & 
\textcolor{white}{\textbf{Energy (kWh)}} \\
\hline
\rowcolor{lightblue}
ResNet18        & 11.7M & 1.8G  & 94.2\% & 0.00093 \\
\hline
MobileNetV2     & 3.5M  & 0.3G  & 92.8\% & 0.00061 \\
\hline
\rowcolor{lightblue}
EfficientNet-B0 & 5.3M  & 0.39G & 93.5\% & 0.00072 \\
\hline
\end{tabular}
\end{center}

\noindent Use \texttt{xcolor} and \texttt{colortbl}.

\vspace{1.5em}

\noindent\textbf{Que 10} Create a student-performance table.

\begin{center}
\begin{tabular}{|l|c|c|c|}
\hline
\rowcolor{headerblue}
\textcolor{white}{\textbf{Student}} & 
\textcolor{white}{\textbf{Mathematics}} & 
\textcolor{white}{\textbf{Physics}} & 
\textcolor{white}{\textbf{Computer Science}} \\
\hline
Amit  & \cellcolor{cellgreen}92  & \cellcolor{cellyellow}81 & \cellcolor{cellgreen}95  \\
\hline
Neha  & \cellcolor{cellred}72    & \cellcolor{cellyellow}88 & \cellcolor{cellgreen}91  \\
\hline
Rahul & \cellcolor{cellred}68    & \cellcolor{cellred}74    & \cellcolor{cellyellow}79 \\
\hline
Priya & \cellcolor{cellgreen}96  & \cellcolor{cellgreen}94  & \cellcolor{cellgreen}98  \\
\hline
\end{tabular}
\end{center}

\noindent\textbf{Que 11} Create a table containing paragraphs inside cells.

\begin{center}
\begin{tabular}{|p{2.3cm}|p{3.8cm}|p{3.5cm}|p{3.2cm}|}
\hline
\textbf{Technique} & \textbf{Description} & \textbf{Advantages} & \textbf{Limitations} \\
\hline
Pruning & Removes unnecessary parameters from a neural network. & Reduces model size and computation. & May reduce accuracy. \\
\hline
Quantization & Represents weights using lower precision. & Reduces memory and energy consumption. & Can introduce numerical error. \\
\hline
Knowledge Distillation & Transfers knowledge from a large teacher model to a smaller student model. & Produces compact models. & Requires training a teacher model. \\
\hline
\end{tabular}
\end{center}

\noindent\textbf{Requirements:}
\begin{itemize}
    \item Use fixed-width columns.
    \item Allow automatic text wrapping.
    \item Properly align text.
    \item Make the table fit within the page.
\end{itemize}

\vspace{2em}

\noindent\textbf{Que 12} Create a professional table using the \texttt{booktabs} package.

\begin{center}
\begin{tabular}{lrrrr}
\toprule
\textbf{Model} & \textbf{Parameters} & \textbf{FLOPs} & \textbf{Accuracy} & \textbf{Energy} \\
\midrule
ResNet18        & 11.7M & 1.82G & 94.21 & 0.00093 \\
MobileNetV2     & 3.5M  & 0.30G & 92.84 & 0.00061 \\
EfficientNet-B0 & 5.3M  & 0.39G & 93.52 & 0.00072 \\
\bottomrule
\end{tabular}
\end{center}

\noindent\textbf{Requirements:}
\begin{itemize}
    \item Use \texttt{\textbackslash toprule}.
    \item Use \texttt{\textbackslash midrule}.
    \item Use \texttt{\textbackslash bottomrule}.
    \item Do not use vertical lines.
\end{itemize}
\noindent\textbf{Que 13} Create a table where one cell contains another table.

\begin{center}
\begin{tabular}{|c|c|c|}
\hline
\textbf{Model} & \textbf{Architecture} & \textbf{Performance} \\ \hline
CNN & $\text{Conv} \rightarrow \text{ReLU} \rightarrow \text{Pool}$ & 
\begin{tabular}{c|c}
Metric & Value \\ \hline
Accuracy & 94\% \\
F1 Score & 93\%
\end{tabular} \\ \hline
ResNet & $\text{Conv} \rightarrow \text{BN} \rightarrow \text{ReLU}$ & 
\begin{tabular}{c|c}
Metric & Value \\ \hline
Accuracy & 96\% \\
F1 Score & 95\%
\end{tabular} \\ \hline
ViT & $\text{Attention} \rightarrow \text{MLP}$ & 
\begin{tabular}{c|c}
Metric & Value \\ \hline
Accuracy & 97\% \\
F1 Score & 96\%
\end{tabular} \\ \hline
\end{tabular}
\end{center}

\vspace{2em}

\noindent\textbf{Que 14} Create the following table containing mathematical expressions.

\begin{center}
\renewcommand{\arraystretch}{2}
\begin{tabular}{|l|c|c|c|}
\hline
\textbf{Model} & \textbf{Loss Function} & \textbf{Parameters} & \textbf{Complexity} \\ \hline
Linear & $\displaystyle L = \frac{1}{n}\sum_{i=1}^{n}(y_i - \hat{y}_i)^2$ & 10K & $O(n)$ \\ \hline
CNN & $\displaystyle L = -\sum y\log(\hat{y})$ & 1.2M & $O(n^2)$ \\ \hline
Transformer & $\displaystyle L = -\sum_i y_i \log p_i$ & 110M & $O(n^2d)$ \\ \hline
\end{tabular}
\end{center}
\end{document}